% Renamed from: "1.3 Information-Conserving Arithmetic and the Dual-Zero Axiom"
% Note (Aug 2026): digit-level m_H=125.1 claims superseded by docs/hardening/ (125 GeV class).
\documentclass[11pt,a4paper]{article}

\usepackage[utf8]{inputenc}
\usepackage{amsmath,amssymb,amsthm}
\usepackage{geometry}
\usepackage{hyperref}
\usepackage{enumitem}

\geometry{margin=1in}

\title{\textbf{Update: Geometric Derivation of the Dual-Zero Regulator}\\
\large SAM3-DualZero-Conoid Framework}
\author{Shawn Dykes (Pieces1121)}
\date{May 2026 (filename cleaned August 2026)}

\begin{document}

\maketitle

\begin{abstract}
We present an updated formulation of the Dual-Zero regulator in the SAM3 spectral triple. The regulator strength $\omega_0$ is derived geometrically from the right conoid structure rather than introduced as a free parameter. This yields $\omega_0 \approx 0.927$. Higgs-sector output is reported in the 125\,GeV class under the August 2026 hardening calibration (digit-level equality with experiment is not claimed here).
\end{abstract}

\section{Information-Conserving Regularization and the Dual-Zero Infinitesimal}
\label{sec:dualzero-axiom}

\subsection{Motivation}
On a non-compact manifold with an infinite discrete spectrum, the high-mode tail of the Dirac operator contributes to the spectral action. To maintain information conservation across scales, we introduce a controlled infinitesimal regularization that never collapses remainders to exact zero while preserving low-energy observables.

\subsection{Definition of the Dual-Zero Algebra}
We work in the truncated hyperreal algebra $\mathbb{R}_{[0,-0]}$ consisting of elements
\[
a + \epsilon b + \epsilon^2 c, \qquad a,b,c \in \mathbb{R},
\]
where the infinitesimal generator is the sequence
\begin{equation}
\epsilon(n) = \omega_0 (-1)^n n^{-n}, \qquad n \in \mathbb{N}, \quad \omega_0 > 0.
\label{eq:epsilon-def}
\end{equation}
Operators acting on the spectrum are regularized symmetrically:
\[
\operatorname{Reg}_2(f)(\lambda) = \lim_{\delta \to 0^+} \frac{f(\lambda + \delta) + f(\lambda - \delta)}{2}.
\]

The regulator strength $\omega_0$ is derived from the geometry of the right conoid:
\begin{equation}
\omega_0 = \left( \frac{R_{\rm curvature}}{D_{\rm bridge}} \right)^{4/13} \approx 0.927.
\label{eq:omega0-derived}
\end{equation}

\subsection{Physical and Geometric Justification}
The form \eqref{eq:epsilon-def} satisfies rapid decoupling, symmetry with the conoid metric, information conservation under $\operatorname{Reg}_2$, and variational consistency of the information current on the critical line (as a selection principle internal to SAM3, not a proof of RH).

\subsection{Impact on Observables}
With geometric $\omega_0$, the spectral action places the Higgs mass in the 125\,GeV class. Corrections to $G_N$ remain negligible at the level discussed in the hardening notes. Newton's constant continues to follow $G_N = 64\pi\ell_0^2/45$, and three chiral generations remain emergent from the $2I$ action.

\subsection{Role in the Framework}
The Dual-Zero algebra is a foundational component of the SAM3 spectral triple, on equal footing with the right conoid geometry and $2I$ symmetry.

\end{document}
